% Compiler: pdfLaTeX
\documentclass[11pt,a4paper]{article}

\usepackage[utf8]{inputenc}
\usepackage[T1]{fontenc}
\usepackage{lmodern}
\usepackage{microtype}
\usepackage{geometry}
\usepackage{enumitem}
\usepackage{titlesec}
\usepackage{fancyhdr}
\usepackage{hyperref}

% Document details
\newcommand{\researchtitle}{Goodhart's law}
\newcommand{\researchshorttitle}{Goodhart's law}
\newcommand{\researchdate}{2026-10-25}

% Layout
\geometry{left=2.7cm,right=2.7cm,top=2.5cm,bottom=2.5cm,headheight=14pt}
\hypersetup{hidelinks,pdftitle={\researchtitle}}
\urlstyle{same}
\frenchspacing
\setlength{\parindent}{0pt}
\setlength{\parskip}{0.55em}
\linespread{1.04}
\setlength{\emergencystretch}{1em}

\titleformat{\section}
    {\large\bfseries}{\thesection}{0.65em}{}
\titleformat{\subsection}
    {\normalsize\bfseries}{\thesubsection}{0.65em}{}
\titlespacing*{\section}{0pt}{1.15em}{0.35em}
\titlespacing*{\subsection}{0pt}{0.7em}{0.2em}
\setlist[itemize]{leftmargin=1.35em,itemsep=0.25em,topsep=0.2em}
\setlist[enumerate]{leftmargin=1.6em,itemsep=0.25em,topsep=0.2em}

\pagestyle{fancy}
\fancyhf{}
\fancyhead[L]{\small Random Research}
\fancyhead[R]{\small\researchshorttitle}
\fancyfoot[C]{\small\thepage}
\renewcommand{\headrulewidth}{0pt}
\renewcommand{\footrulewidth}{0pt}

\begin{document}

{\LARGE\bfseries\researchtitle\par}
\vspace{0.4em}
{\small\researchdate\par}

\section{Overview}

Goodhart's law concerns what happens when we push a measure hard enough
that it stops tracking the thing we care about. A count can help us
understand an activity, then become a target that changes the activity
itself. \cite{manheim}

For a beginner, the useful question is simple: could this number improve
while the actual outcome gets worse?

\section{Findings}

Strathern examines assessment in British universities. Her account treats
audit as something that changes institutional life, with consequences
beyond recording what was already happening. That helps explain why
measurement deserves attention as part of a system of incentives.
\cite{strathern}

Manheim and Garrabrant distinguish several mechanisms behind failures of
measurement. A measure can contain noise, or a relationship can stop
holding when pushed beyond the conditions in which it was useful.
Their account also separates those problems from deliberate manipulation.
Cheating is therefore only one possible explanation. \cite{manheim}

The paper is a conceptual analysis. It offers ways to examine a failure;
it does not provide a universal threshold at which any particular metric
will stop working.

\section{Example}

Suppose a support team is rewarded for closing tickets quickly. An
employee could raise the closure count by ending conversations before
the problem is resolved. Another could spend time on a difficult case
and close fewer tickets despite helping the customer.

This hypothetical example does not show that speed is a bad measure.
It shows why a closure count cannot settle the question of whether the
customer received useful help. Checking repeat contacts or reading a
sample of conversations could reveal something the count leaves out.

\section{Discussion}

I like this idea because it makes a familiar frustration easier to
describe. People can follow the target they were given and still produce
work that disappoints whoever set it.

My hesitation is that Goodhart's law can become a convenient excuse to
reject measurement altogether. Without any record, judging performance
can become arbitrary too. I would keep a useful measure and be willing
to question what it leaves out.

I would also avoid assuming bad faith. An employee may be adapting
sensibly to the incentives of a job, while the organization has defined
success too narrowly.

\section{Practical implications}

Before using a target, I would write down the outcome it is meant to
represent. Then I would try to construct a case in which the target
improves without that outcome improving.

That case suggests what else to inspect. A second measure might help,
although it can develop the same problem. Occasional review of the
underlying work still seems necessary.

\renewcommand{\refname}{References}
\begin{thebibliography}{9}
\small
\setlength{\itemsep}{0.35em}

\bibitem{strathern}
Strathern, M. (1997).
\href{https://www.cambridge.org/core/journals/european-review/article/improving-ratings-audit-in-the-british-university-system/FC2EE640C0C44E3DB87C29FB666E9AAB}
{Improving ratings: Audit in the British University system.}
\textit{European Review}, 5(3), 305--321.

\bibitem{manheim}
Manheim, D., and Garrabrant, S. (2019).
\href{https://arxiv.org/abs/1803.04585}
{Categorizing variants of Goodhart's law.}
arXiv:1803.04585, version 4. First submitted in 2018.

\end{thebibliography}
\end{document}
